% GENERATED FILE. Edit canonical lessons, then run npm run generate:books.
% canonical-source-hash: fnv1a64:91bebf8d6e9aa844
% canonical-lessons: ZH-C178-zhengchang, ZH-W178-chang-often, ZH-C178-rizi, ZH-C178-congqian, ZH-C178-zuijin, ZH-C178-lingchen

\chapter{An Ordinary Day}
\label{ch:zh-an-ordinary-day}
\hlchaptermodality{178}

\begin{chapteropening}
\textbf{By the end of this chapter:} \emph{I can say what is normal, talk about my days, then and lately, and the small hours.}

Complete the last lesson of chapter 178: I can say what is normal, talk about my days, then and lately, and the small hours.
\end{chapteropening}

\section[\texorpdfstring{\zh{正常} (zhèngcháng)}{正常 (zhèngcháng)}]{\zh{正常} (zhèngcháng) — normal; as usual}

\label{lesson:ZH-C178-zhengchang}



\begin{quote}
Before the new one: say the Mandarin for fresh, then the Mandarin for horrible to eat.
\end{quote}

\subsection*{You'll want to know: \zh{正常}}
\textbf{\zh{正常}} — \emph{zhèngcháng} — \textquotedblleft{}normal; as usual\textquotedblright{}.

\begin{quote}
\textbf{\zh{没事我觉得很正常}} \emph{méishì wǒ juéde hěn zhèngcháng} — It's nothing: I think that's perfectly normal.
\end{quote}

The new character, in pieces: \zh{常} — eleven strokes: three short strokes on top, the cover piece of \zh{学}, \zh{口} under it, and at the foot a frame with a long vertical down through it. The next lesson puts it in your hand.

\begin{cousinweb}
\textbf{\zh{正常}} (\emph{zhèngcháng}) is built from \zh{正} 'right, proper' and \zh{常}, both inherited from Middle Chinese, \emph{tsyengH} and \emph{dzyang} (Guangyun, in Baxter's transcription). The dictionary glosses the second as a length of two fathoms and as an official's title; the sense 'usual' is not given there, so how the pair came to mean normal is not settled.
\end{cousinweb}

\begin{grammarlens}[title={What you've built}]
One more word for times of day.
\end{grammarlens}

\subsection*{Your turn}
\begin{itemize}
\raggedright
  \item \emph{Say it:} \emph{zhèngcháng}
  \item \emph{Say it:} \emph{zhèngcháng}, once more
  \item \emph{Say it:} \emph{méishì wǒ juéde hěn zhèngcháng} — It's nothing: I think that's perfectly normal.
  \item \emph{Recall:} say the Mandarin for fresh, then the Mandarin for horrible to eat, then say \emph{zhèngcháng} again
\end{itemize}

\begin{tcolorbox}[breakable,colback=teal!4,colframe=teal!35!black,title={Before you move on}]
What does \zh{正常} mean? (\textquotedblleft{}Normal; as usual\textquotedblright{}.) Say it once more.
\end{tcolorbox}
\section[\texorpdfstring{\zh{常} (cháng)}{常 (cháng)}]{\zh{常} — three strokes, a cover, \zh{口} and a cloth}

\label{lesson:ZH-W178-chang-often}



\begin{quote}
The word first, then the character inside it.

\begin{itemize}
\raggedright
  \item \emph{Recall:} say \emph{zhèngcháng} — \textbf{R1}, one lesson back, before it decays
\end{itemize}

\begin{quote}
\zh{学} \zh{口}
\end{quote}

\emph{Read it:} each one aloud, with nothing said first — \emph{xué}, \emph{kǒu}
\end{quote}

\subsection*{Script — three strokes, a cover, \zh{口} and a cloth}
\hlblockfigure{\detokenize{figures/ZH-W178-chang-often-filmstrip.pdf}}{How \zh{常} is written, stroke by stroke}

\begin{quote}
\zh{常}
\end{quote}

\textbf{Four pieces, top to bottom.} Count the pieces, not the strokes.

\textbf{Three short strokes on top:} a vertical in the middle, a dot to its left, and a short stroke on its right falling to the left.

\textbf{Then the cover piece of \zh{学}, two strokes:} a short stroke falling down on the left, then a horizontal that turns down and hooks.

\textbf{Then \zh{口}, three strokes,} as you write it.

\textbf{At the foot, the cloth piece, three strokes:} a short vertical on the left, a horizontal that turns down and hooks, and last a long vertical down through the middle.

3 + 2 + 3 + 3 = eleven.

\subsection*{Your turn}
\begin{itemize}
\raggedright
  \item \emph{Write it:} \zh{常}, all eleven, with the model in front of you
  \item \emph{Cover:} write it again from memory, and count the pieces, not the strokes
\end{itemize}

\begin{tcolorbox}[breakable,colback=teal!4,colframe=teal!35!black,title={Before you move on}]
How many strokes in \zh{常}? (\textbf{Eleven}.) What sits under the cover? (\textbf{\zh{口}}.) Which stroke comes last? (\textbf{The long vertical through the cloth piece}.)

Source: \href{https://raw.githubusercontent.com/chanind/hanzi-writer-data/68d10a4b21150cae5e1ebbd223eed289cf32d90c/data/\%E5\%B8\%B8.json}{Hanzi Writer Data \zh{常}, pinned stroke paths and medians}.
\end{tcolorbox}
\section[\texorpdfstring{\zh{日子} (rìzi)}{日子 (rìzi)}]{\zh{日子} (rìzi) — day; days (of one's life)}

\label{lesson:ZH-C178-rizi}



\begin{quote}
Before the new one: say the Mandarin for horrible to eat, then the Mandarin for normal.
\end{quote}

\subsection*{You'll want to know: \zh{日子}}
\textbf{\zh{日子}} — \emph{rìzi} — \textquotedblleft{}day; days (of one's life)\textquotedblright{}.

\begin{quote}
\textbf{\zh{小时候的日子很快乐}} \emph{xiǎoshíhou de rìzi hěn kuàilè} — My childhood days were happy.
\end{quote}

\begin{cousinweb}
\textbf{\zh{日子}} (\emph{rìzi}) is built from \zh{日} 'sun, day' and the noun ending \zh{子}, both inherited from Middle Chinese, \emph{nyit} and \emph{tsiX} (Guangyun, in Baxter's transcription). In the word the second loses its tone.
\end{cousinweb}

\begin{grammarlens}[title={What you've built}]
One more word for times of day.
\end{grammarlens}

\subsection*{Your turn}
\begin{itemize}
\raggedright
  \item \emph{Say it:} \emph{rìzi}
  \item \emph{Say it:} \emph{rìzi}, once more
  \item \emph{Say it:} \emph{xiǎoshíhou de rìzi hěn kuàilè} — My childhood days were happy.
  \item \emph{Recall:} say the Mandarin for horrible to eat, then the Mandarin for normal, then say \emph{rìzi} again
\end{itemize}

\begin{tcolorbox}[breakable,colback=teal!4,colframe=teal!35!black,title={Before you move on}]
What does \zh{日子} mean? (\textquotedblleft{}Day; days (of one's life)\textquotedblright{}.) Say it once more.
\end{tcolorbox}
\section[\texorpdfstring{\zh{从前} (cóngqián)}{从前 (cóngqián)}]{\zh{从前} (cóngqián) — once upon a time; in the past}

\label{lesson:ZH-C178-congqian}



\begin{quote}
Before the new one: say the Mandarin for normal, then the Mandarin for one's days.
\end{quote}

\subsection*{You'll want to know: \zh{从前}}
\textbf{\zh{从前}} — \emph{cóngqián} — \textquotedblleft{}once upon a time; in the past\textquotedblright{}.

\begin{quote}
\textbf{\zh{我从前常常喝茶}} \emph{wǒ cóngqián chángcháng hē chá} — I used to drink a lot of tea.
\end{quote}

\begin{cousinweb}
\textbf{\zh{从前}} (\emph{cóngqián}) is built from \zh{从} 'from' and \zh{前} 'before', both inherited from Middle Chinese, \emph{dzjowng} and \emph{dzen} (Guangyun, in Baxter's transcription): from before.
\end{cousinweb}

\begin{grammarlens}[title={What you've built}]
One more word for times of day.
\end{grammarlens}

\subsection*{Your turn}
\begin{itemize}
\raggedright
  \item \emph{Say it:} \emph{cóngqián}
  \item \emph{Say it:} \emph{cóngqián}, once more
  \item \emph{Say it:} \emph{wǒ cóngqián chángcháng hē chá} — I used to drink a lot of tea.
  \item \emph{Recall:} say the Mandarin for normal, then the Mandarin for one's days, then say \emph{cóngqián} again
\end{itemize}

\begin{tcolorbox}[breakable,colback=teal!4,colframe=teal!35!black,title={Before you move on}]
What does \zh{从前} mean? (\textquotedblleft{}Once upon a time; in the past\textquotedblright{}.) Say it once more.
\end{tcolorbox}
\section[zuìjìn]{zuìjìn — lately; recently}

\label{lesson:ZH-C178-zuijin}



\begin{quote}
Before the new one: say the Mandarin for one's days, then the Mandarin for once upon a time.
\end{quote}

\subsection*{You'll want to know: zuìjìn}
\emph{zuìjìn} — \textquotedblleft{}lately; recently\textquotedblright{}.

\begin{quote}
\emph{nǐ zuìjìn máng ma} — Have you been busy lately?
\end{quote}

Its characters come later; for now it is a word for the ear.

\begin{cousinweb}
\emph{zuìjìn} is built from \emph{zuì} 'most' and \emph{jìn} 'near', both inherited from Middle Chinese, \emph{tswajH} and \emph{gj+nX} (Guangyun, in Baxter's transcription): the nearest time.
\end{cousinweb}

\begin{grammarlens}[title={What you've built}]
One more word for times of day.
\end{grammarlens}

\subsection*{Your turn}
\begin{itemize}
\raggedright
  \item \emph{Say it:} \emph{zuìjìn}
  \item \emph{Say it:} \emph{zuìjìn}, once more
  \item \emph{Say it:} \emph{nǐ zuìjìn máng ma} — Have you been busy lately?
  \item \emph{Recall:} say the Mandarin for one's days, then the Mandarin for once upon a time, then say \emph{zuìjìn} again
\end{itemize}

\begin{tcolorbox}[breakable,colback=teal!4,colframe=teal!35!black,title={Before you move on}]
What does zuìjìn mean? (\textquotedblleft{}Lately; recently\textquotedblright{}.) Say it once more.
\end{tcolorbox}
\section[língchén]{língchén — the small hours; before dawn}

\label{lesson:ZH-C178-lingchen}



\begin{quote}
Before the new one: say the Mandarin for once upon a time, then the Mandarin for lately.
\end{quote}

\subsection*{You'll want to know: língchén}
\emph{língchén} — \textquotedblleft{}the small hours; before dawn\textquotedblright{}.

\begin{quote}
\emph{wǒ língchén sì diǎn qǐchuáng le} — I got up at four in the morning.
\end{quote}

Its characters come later; for now it is a word for the ear.

\begin{cousinweb}
\emph{língchén} is built from \emph{líng} and \emph{chén}, both inherited from Middle Chinese, \emph{ling} and \emph{dzyin} (Guangyun, in Baxter's transcription). The dictionary glosses the first 'ice' and the second 'early' and 'dawn'; how the first came into the word is not settled.
\end{cousinweb}

\begin{grammarlens}[title={What you've built}]
One more word for times of day.
\end{grammarlens}

\subsection*{Your turn}
\begin{itemize}
\raggedright
  \item \emph{Say it:} \emph{língchén}
  \item \emph{Say it:} \emph{língchén}, once more
  \item \emph{Say it:} \emph{wǒ língchén sì diǎn qǐchuáng le} — I got up at four in the morning.
  \item \emph{Recall:} say the Mandarin for once upon a time, then the Mandarin for lately, then say \emph{língchén} again
\end{itemize}

\begin{tcolorbox}[breakable,colback=teal!4,colframe=teal!35!black,title={Before you move on}]
What does língchén mean? (\textquotedblleft{}The small hours; before dawn\textquotedblright{}.) Say it once more.
\end{tcolorbox}
